\section{Literature Review}
\label{sec:related-work}
To ground our benchmark design in existing research and ensure a realistic sequence generator, 
this section reviews prior benchmarking approaches for client-side query optimization and real-world query usage patterns.

\subsection{Existing Benchmarking Approaches}
For the evaluation of client-side benchmarking techniques that rely on previous queries,
we primarily investigate caching literature, as this is a well-studied problem relevant to our work.

In the literature, we find three distinct benchmarking approaches: 
simulated micro benchmarks, simulated end-to-end query benchmarks, and real-world query logs.

% Existing evaluation methods for client-side query optimization 
% generally fall into three categories:


\textbf{Simulated micro benchmarks} attempt to isolate performance gains of the benchmarked approaches by fixing certain aspects of query execution.
This can range from fixing the cached data size~\cite{hartig2011caching} to controlling the percentage of query-relevant data present in the cache \cite{nicholson2024hpcache}. Fixing aspects of query execution makes it easier to isolate and understand the effect of specific caching mechanisms. 
However, it also gives an incomplete picture, because we don’t know how those fixed conditions relate to real workloads or how the methods perform when those assumptions don’t hold.

\textbf{Simulated end-to-end query benchmarks} use generated query sequences to evaluate performance. 
They aim to measure cache hit rate and overall execution behavior under the assumption that the simulated sequences reflect realistic usage.
In practice, the simulation is usually built by extending an existing benchmark so that it generates sequences instead of isolated queries. 
The sequence order is often determined by simple rules or by random selection.

A common approach is to use queries from existing benchmarks such as the Star Schema Benchmark \cite{o2009star} or the Berlin SPARQL Benchmark~\cite{bizer2011berlin}.
These queries are then randomly divided into sequences \cite{hartig2011caching,nicholson2024hpcache,folz2016cyclades}, with each sequence executed using a shared cache. This approach is straightforward but not realistic.
Prior work shows that real workloads contain user behavior patterns and temporal locality, and these patterns can be exploited by caching methods.
Random sequences fail to capture these properties.

More realistic is to adjust the query generators of existing benchmarks to include patterns in the query sequences.
Some works~\cite{williams2011enabling,martin2010improving} use a Pareto distribution to decide what entity is used to instantiate a query template.
This follows from the observation that many human activities follow a Pareto distribution~\cite{arnold2014pareto},
 with many queries relating to a small portion of the entities.

A third approach uses simulated query sequences that reflect the access patterns of specific applications.
One example is a workload that models data scientists exploring different aspects of the same entity—such as an author in a scholarly network—through consecutive, thematically related queries~\cite{chatzopoulos2022atrapos}.
Such simulations often include multiple query sessions, with new sessions beginning according to a probability parameter $p$.
Another example is a simulated application scenario focused on retrieving information about vacation destinations, which is used to assess caching effectiveness~\cite{martin2010improving}.

\textbf{Real-world query logs} Whenever available, query logs from actual applications form a highly relevant benchmarking approach.
However, obtaining these query logs can be challenging.
For SPARQL query optimization, query logs for a variety of endpoints are available, such as DBpedia~\cite{saleem2015lsq,stadler2024lsq,berendt2013usewod2013} and Wikidata~\cite{stadler2024lsq}.
These logs are extensively used in benchmarking SPARQL-based approaches~\cite{zhang2016secf,papailiou2015graph,lorey2013caching}.

The evaluation of personalized LTQP optimization would ideally use real-world query logs, but such logs do not exist.
To avoid the fragmented evaluation practices of current simulated approaches, the next best option is a unified simulated benchmark. 
A shared benchmark would provide a common basis for comparison, improve reproducibility, and prevent each paper from recreating its own evaluation setup.

\subsection{Real-world Query Usage Patterns}
\label{sec:related_work_qup}
To inform the benchmark design on what patterns should be simulated, we will outline three relevant client query usage patterns observed in the literature. The relevance is determined for our social network application use case.

\textbf{User Interest-Based Query Behavior}
Work on social networks shows users form communities based on shared interests~\cite{mislove2007measurement,clark2007understanding,backstrom2006group,kalaitzakis2012evolution}. 
These networks exhibit power-law degree distributions, small-world structures, and scale-free properties. 
This topology creates a dense core of high-degree nodes that connect many smaller, strongly clustered groups~\cite{mislove2007measurement}. 
Simulation studies confirm that this interest-driven group formation is stable and reliably modeled~\cite{amblard2015models,hamill2009social}.

Furthermore, empirical research formally establishes that shared interests directly drive a higher degree of explicit and implicit interactions between users~\cite{mitzlaff2013user}. 
These findings justify our simulation of user-relatedness and interest-driven behavior when generating query sequences.

\textbf{Logical Query Sessions}
Analysis of web browsing behavior shows that activity is best grouped into \emph{logical sessions}: sequences of user interactions focused on specific goals, 
rather than purely time-bound intervals~\cite{meiss_whats_2009,sadagopan2008characterizing,schneider2009understanding}. 
These task-oriented sessions frequently involve non-linear navigation, 
including backtracking. 
New sessions begin in approximately $15\%$ of clicks, 
following a log-normal distribution with average size and depth means of 6.1 and 3, respectively. 

These patterns also appear in application navigation. 
In a decentralized social media application, for example, 
a user viewing a feed might click a creator's username, view comments, or like a post. 
Because the parameters of each new query derive directly from the preceding query's results, 
these queries chain together into structures mirroring logical web sessions. 
Additionally, direct searches for specific posts or topics reflect standard session-switching behavior. 

Consequently, interactive applications generate sequences of interdependent requests~\cite{vogele2018wessbas}. 
This dynamic informs the design of Facebook's TAO data store, 
which is explicitly optimized for workloads driven by the user's step-by-step traversal of the social graph within the interface~\cite{bronson2013tao}.

\textbf{Query Refinement}  
Analysis of SPARQL queries issued to the DBpedia endpoint reveals that 
users frequently submit related queries in streaks~\cite{bonifati2020analytical,zhang2020characterizing}. 
These streaks consist of query sequences that share an underlying intention but undergo 
iterative refinement to achieve the user's objective. 
Previous studies initially introduced the concept of SPARQL query sessions specifically to analyze 
robotic queries~\cite{zhang2020characterizing,bonifati2020analytical}.

In this benchmark, we focus on organic query sequences~\cite{zhang2020characterizing}, 
which users generate through interactions with software or browsers~\cite{malyshev2018getting}.
Because organic interactions are driven by application design rather than underlying data organization, 
we assume these user behaviors will also emerge in applications navigating structured decentralized environments
Furthermore, we deprioritize robotic queries because their sequences can 
already be simulated by traditional benchmark approaches~\cite{taelman2023link,bizer2011berlin,alucc2014diversified}.

For both organic and robotic queries, analysis of total usage and reformulations (removals and additions) 
of operators over all query logs shows that filter, union, and optional operators account for the vast majority of reformulated queries.

Recent analyses of Wikidata query logs further demonstrate that query development is an iterative process of design,
debugging, and maintenance~\cite{arenas2025exploring}.
Based on this empirical data, 
exploratory querying encompasses six core behaviors~\cite{arenas2025exploring}:
\textbf{result refinement} to reduce records via filters or selective triple patterns,
\textbf{result generalization} to retrieve more entities by broadening parameters,
such as adding superclasses or increasing limits,
\textbf{result extension} to fetch additional attributes by appending new triple patterns,
\textbf{bug fixing} to correct structural errors like invalid URIs,
\textbf{query refinement} to adjust complex language features,
and \textbf{shifting query intent} to pivot topics by altering core predicates.

These findings confirm that real-world querying relies on dynamic, step-by-step modifications rather than the static execution of independent templates.
Consequently, existing evaluation frameworks fail to capture how users actively interact with data systems. 
Evaluating system performance for these exploratory query sequences remains a critical research challenge, 
explicitly requiring new benchmarks to validate progress~\cite{arenas2025exploring}. 
Therefore, modeling these empirical sequence patterns is essential to accurately assess client-side 
query optimization techniques.